% Generated by tools/paper/render.py from main_text.py; edit the source there, not this file.
\section{Definitions}

Sequences are strings over \(\{A, C, G, T\}\); any other symbol, such as N, matches nothing, so no k-mer or match spans it. The reference \(R\) may contain several sequences, and matches do not cross sequence boundaries. \(S[i..j)\) denotes the substring of \(S\) from position \(i\) to \(j - 1\), and \(\overline{S}\) its reverse complement.

\textbf{Sampled k-mers.} For a hash \(h\) on k-mers, the (w, k)-minimizer of a window of w consecutive k-mers is its k-mer of minimum hash, with ties broken towards the leftmost; the minimizers of a sequence are those of all its windows, each reported once \citep{Roberts2004}. Canonical minimizers use a hash that is invariant under reverse complementation. A k-mer is a \emph{closed syncmer} if its minimal s-mer starts at offset 0 or \(k - s\), and an \emph{open syncmer} with offset t if it starts at t \citep{Edgar2021}. For comparison with simd-minimizers we use its parameterization, in which a syncmer is defined by the length k of the ranked substrings and their number w, and spans \(k + w - 1\) bases \citep{GrootKoerkamp2025simd}.

\textbf{Exact matches between genomes.} A \emph{MEM} of a query sequence \(Q\) and \(R\) is a triple \((r, q, \ell)\) with \(R[r..r+\ell) = Q[q..q+\ell)\) and \(\ell \ge L\) that cannot be extended to the left or right. It is a \emph{MUM} if its string occurs exactly once in \(R\) and once in \(Q\), and \emph{reference-unique} (MUMmer's \texttt{-mumreference}) if it occurs once in \(R\) \citep{Kurtz2004}. It is \emph{frequency-capped} at N if its string occurs at most N times in \(R\) or, in the strict form, at most N times in both \(R\) and \(Q\); a floor M additionally requires at least M occurrences in \(R\). N = 1 yields the MUMs. Reverse-strand matches are those of \(\overline{Q}\), with positions along \(\overline{Q}\), as in MUMmer.

\textbf{Exact matches of reads.} An interval \(Q[i..j)\) of a read \emph{occurs} if it or its reverse complement is a substring of \(R\). An \emph{SMEM} is an interval of length at least L that occurs and is not contained in a longer occurring interval \citep{Li2012}. An \emph{f-SMEM} replaces ``occurs'' by ``occurs at least f times'' \citep{Tatarnikov2022,Gagie2026}; f = 1 gives the SMEMs, and a cap N removes SMEMs that occur more than N times.
